\documentclass[instructions.tex]{subfiles}
\begin{document}
 
\chapter{Tout passe ici bas. Réflexions, etc.}
 

\primo Ce qui se passe doit être estimé comme rien\footnote{Quod aeternum non est, nibil est}. Cent années sont donc peu de chose, puisqu'elles s'écoulent à tout instant et qu'elles passent; elles sont moins qu'un moment, en comparaison de l'avenir. 

Imaginez-vous, d'un côté, une vie austère et pénitente, une maladie longue, une disgrâce accablante; tout cela sera, dans quelques temps, aussi véritablement passé que s'il n'avait jamais été; ce qui en reste, c'est l'espérance d'une récompense immortelle. Voilà ce qui console un chrétien dans les tribulations de cette vie. 

Imaginez-vous, d'un autre côté, de grandes richesses avec tous les plaisirs de cette vie: un jour viendra qu'il sera vrai de dire que tout cela n'est plus. Tout ce qui en restera, sera le chagrin de tout quitter en un moment, avec la crainte d'un châtiment éternel; et voilà ce qui afflige et accable un mondain. 

Que reste-t-il de la puissance d'Alexandre, des richesses de Crésus, de la beauté de Cléopâtre? on les admirait pendant leur vie, on en parle après leur mort, tandis qu'ils sont malheureux. On les louent où ils ne sont plus, dit saint Augustin, et on les tourmente où ils sont\footnote{Lauduntur ubi non cunt, cruciantur ubi sunt. S. Aug.}. Le temps du salut est écoulé pour eux; il ne l'est point encore pour vous, mais il le sera bientôt\dots

Il y a cent ans que vous n'étiez-pas; avant que cent ans soient finis, vous ne serez plus: le temps du salut sera passé, votre corps réduit en poussière, et vote âme jugée. Est-il possible que cette réflexion ne puisse dessiller vos yeux, et détacher votre coeur des biens et des plaisirs de la vie, que vous quitterez dans peu de jours, quelques effort que vous fassiez pour les retenir ?

Un grand roi disait, à sa mort:\og Tout est passé pour moi. Ah! qu'il me serait bien plus avantageux d'avoir été pauvre villageois craignant Dieu, que d'avoir été un puissant monarque!\fg Que nous sommes aveugles de nous laisser tromper par des choses qui n'ont que l'apparence! Rien ne mérite ici-bas nos empressemens, parce que rien n'y est durable. Tout ce que le temps peut nous ravir, dis saint Eucher, est peu de chose et n'est pas digne de nous\footnote{Nihil est magnum re, quod parvum tempore. S.Euch}. Comblez un homme de tout ce que son coeur peut désirer, il n'est plus heureux, dés qu'il sait que bientôt tout sera passé pour lui. 

Cette pensée:\og Tout passe\fg, est comme un fiel qui devrait nous ôter le goût des choses d'ici-bas. \og Combien durera cette fortune, ce plaisir, cet honneur, cette beauté?\fg oh! que cette réflexion a sanctifié de personnes en leur découvrant, même à la fleur de leur âge, la vanité de tout ce que le monde a de plus pompeux et de plus florissant! \og Ne me parlez point\fg, disait un pieux cardinal en mourant, \og ne me parlez point de toutes les grandeurs humaines. Je m'estimerais bien plus heureux d'avoir été le dernier dans un cloitre, que de me voir revêtu de la pourpre\fg. 

Les plus grandes afflictions commencent souvent dans les plus grands biens et par les plus grands plaisirs; et ces afflictions sont d'autant plus sensibles, que ces biens et ces plaisirs finissent lorsque nous y pensons le moins. Ne vaut-il pas mieux s'en détacher à présent avec mérite, que de les perdre un jours avec douleur et sans fruit? 

\secundo Les plaisirs de la vie, fussent-ils solides, ce serait toujours une folie de s'y attacher. Ils répandent des ténèbres dans notre esprit et nous font oublier l'autre vie, pour laquelle seule nous devons travailler. En effet, puisque la félicité et les joies du ciel ne passeront jamais, pourquoi ne vous attachez-vous pas à ce qui peut vous les procurer? Un homme qui doit voyager pendant une année sur la mer, ne serait-il pas insensé si, pour y vivre, il faisait provision d'alimens qui se corrompraient dans un jour ? Et vous qui voyagez pour l'éternité, n'êtes vous pas encore plus aveugle, si, au lieu d'acquérir un fonds de mérites pour y vivre heureux, vous ne pensez qu'à vous procurer des plaisirs et des biens caducs qui, dans deux jours, vous seront-inutiles ? 

Cette attache aux biens de la terre est, au sentiment même des païens, la marque d'un petit esprit et d'une âme bornée dans ses sentimens\footnote{Angustus est animus, quem terrena delectant. Senec.}. Tout ce que vous pouvez espérer ou posséder ici-bas doit périr dans peu de temps; et, tout étant péri pour vous, que deviendrez-vous? Aimez donc et cherchez ce qu'on peut posséder sans crainte de le perdre, qui est le ciel et Dieu seul. Dans tout le reste on n'éprouve qu'inconstance et disgrâces: souvent le soir accable de tristesse ceux que le matin avait comblés de joie; c'est à dire qu'une vie qu'on a commencée et passée dans le plaisir et l'abondance, finit par le regret.

\section{Résolutions}

\primo Je travaillerais donc à me détacher du monde, de ses plaisirs, de ses honneurs et de ses biens périssables.

\secundo Et je porterai toutes mes affections dans le ciel, où sont les biens véritables et permanens.

Mon Dieu, aidez-moi à sortir en esprit de cette terre pleine d'illusions et d'erreurs; où tout passe, et à élever mes pensées et mon coeur dans le séjour de la gloire et des délices éternelles. 

\end{document}
